\chapter[Mesh Independence of Validation Geometries]{Mesh Independence of the Extracted Validation Geometries}
\label{app:grid_independence_optimized}

\providecommand{\meshplaceholder}[1]{%
	\fbox{\parbox[c][0.20\textheight][c]{0.88\linewidth}{\centering \textbf{Placeholder}\\[0.5em]#1}}%
}

This appendix supports the Ansys re-analysis campaign described in Chapter~\ref{ch:benchmark_results}. Once the topology optimization produces a final projected design field, the geometry must be extracted, cleaned, and meshed again using a body-fitted CFD mesh. The resulting pressure-drop values are only meaningful if this remeshing step is itself spatially converged.

\section{Mesh Generation for Extracted Geometries}
The extracted validation geometries differ from the native FEniCS optimization domain in two important ways. First, the Brinkman-penalized diffuse interface is replaced by a sharp wall. Second, the validation mesh must now satisfy the near-wall requirements of the selected Ansys turbulence model rather than those of the topology optimization mesh.

For SA and SST $k$-$\omega$, the preferred strategy remains a wall-resolved mesh. For standard $k$-$\epsilon$, the near-wall treatment must be stated explicitly in the final manuscript, because the required mesh structure depends on whether wall functions or a more resolved near-wall treatment is used. In all cases, the extracted contour threshold, the amount of geometric smoothing, and the treatment of non-design inlet and outlet regions should be documented alongside the mesh statistics.

\section{Mesh Study Template by Benchmark}
Table~\ref{tab:app_validation_mesh_template} is intended as the minimum record for the Ansys mesh-independence study. It should be completed separately for each geometry and, where the wall treatment differs, separately for each turbulence model.

\begin{table}[htbp]
	\centering
	\footnotesize
	\caption{Template for the Ansys mesh-independence study on extracted geometries.}
	\label{tab:app_validation_mesh_template}
	\renewcommand{\arraystretch}{1.2}
	\begin{tabularx}{\textwidth}{l l c c c X}
		\hline
		\textbf{Case} & \textbf{Model} & \textbf{Coarse} & \textbf{Medium} & \textbf{Fine} & \textbf{Comment} \\
		\hline
		Diffuser & SA & [TBD] & [TBD] & [TBD] & Report cells, first-layer height, and metric asymptote. \\
		Diffuser & SST $k$-$\omega$ & [TBD] & [TBD] & [TBD] & State whether the same mesh can be reused. \\
		Double pipe & SA & [TBD] & [TBD] & [TBD] & Include junction refinement strategy. \\
		Double pipe & SST $k$-$\omega$ & [TBD] & [TBD] & [TBD] & Report sensitivity of central mixing region. \\
		Curved benchmark & SA & [TBD] & [TBD] & [TBD] & Include near-wall curvature resolution. \\
		Curved benchmark & SST $k$-$\omega$ & [TBD] & [TBD] & [TBD] & Report whether separation location is mesh-stable. \\
		Curved benchmark & Standard $k$-$\epsilon$ & [TBD] & [TBD] & [TBD] & Only if this model is retained in the final campaign. \\
		\hline
	\end{tabularx}
\end{table}

\section{Convergence of Performance Metrics}
The preferred output of the mesh study is not only the raw cell count, but the asymptotic behavior of the chosen comparison metric. This should be the same metric used in Chapter~\ref{ch:benchmark_results}, preferably pressure drop and optionally hydraulic power loss for fixed-flow-rate cases.

\begin{figure}[htpb]
	\centering
	\begin{subfigure}[b]{0.48\textwidth}
		\centering
		\meshplaceholder{Diffuser mesh-independence plot: comparison metric versus mesh level.}
		\caption{Diffuser}
		\label{fig:app_mesh_diffuser}
	\end{subfigure}
	\hfill
	\begin{subfigure}[b]{0.48\textwidth}
		\centering
		\meshplaceholder{Double-pipe mesh-independence plot: comparison metric versus mesh level.}
		\caption{Double pipe}
		\label{fig:app_mesh_doublepipe}
	\end{subfigure}

	\vspace{0.5cm}

	\begin{subfigure}[b]{0.48\textwidth}
		\centering
		\meshplaceholder{Curved-benchmark mesh-independence plot: comparison metric versus mesh level.}
		\caption{Curved benchmark}
		\label{fig:app_mesh_curved}
	\end{subfigure}
	\hfill
	\begin{subfigure}[b]{0.48\textwidth}
		\centering
		\meshplaceholder{Optional plot of local profile convergence at the most sensitive section.}
		\caption{Profile-based convergence check}
		\label{fig:app_mesh_profiles}
	\end{subfigure}
	\caption{Placeholder figure set for the mesh-independence evidence required by the Chapter~\ref{ch:benchmark_results} Ansys campaign.}
	\label{fig:app_opt_mesh_study}
\end{figure}

In the final text, the interpretation should remain modest. The goal is not to claim a perfect mesh, but to show that the comparison metric has stabilized enough that the differences reported between FEniCS-SA, Ansys-SA, and Ansys SST $k$-$\omega$ are dominated by the solver or turbulence model rather than by remeshing noise.
